% !TEX program = xelatex
\documentclass[11pt,a4paper]{ctexart}

\usepackage{amsmath,amssymb,amsfonts}
\usepackage{booktabs}
\usepackage{geometry}
\usepackage{hyperref}
\usepackage{enumitem}
\usepackage{xcolor}
\usepackage{longtable}
\usepackage{array}

\geometry{left=2.2cm,right=2.2cm,top=2.4cm,bottom=2.4cm}
\hypersetup{colorlinks=true,linkcolor=blue,urlcolor=blue}

\newcolumntype{L}[1]{>{\raggedright\arraybackslash}p{#1}}
\newcolumntype{C}[1]{>{\centering\arraybackslash}p{#1}}
\newcolumntype{T}[1]{>{\ttfamily\raggedright\arraybackslash\footnotesize}p{#1}}
\newcommand{\tid}[1]{\texttt{\detokenize{#1}}}

\title{%
  \textbf{星闪 SLB 物理层}\\[0.35em]
  \Large L0 + L1 参数登记册（初稿）\\[0.25em]
  \large 全 PHY 共用常量 \textbullet\ 上层下发会话配置 \textbullet\ 读写所有权
}
\author{依据当前工作区技术文档分册整理（6.2 / 6.5 / 6.6）\\
精确数值与字段定义以 T/XS 10001-2025 正式标准为准\\[0.3em]
\small 工程用途：AI 分模块实现时的跨模块参数真源（SSOT）初稿；\\
\small 命名约定 airMode\,=\,\texttt{slb}（nearlink-naming）
}
\date{2026 年 7 月 27 日}

\begin{document}
\maketitle
\tableofcontents
\newpage

%======================================================================
\part{总览}
%======================================================================

\section{文档目的与范围}

本文档汇总 SLB 120\,kHz 物理层实现中应统一管理的两类参数：

\begin{enumerate}[nosep]
  \item \textbf{L0 系统常量}：全 PHY 共用、基本不变的时频/网格/物理常数；
  \item \textbf{L1 会话配置}：由上层（MAC / 第 7 章解析 / 调度）下发，PHY 各模块只读消费的量。
\end{enumerate}

\textbf{不在本册展开}：L2（模块专有可调参数）、L3（运行时测量/命令交换量）。二者见 README 后续安排，随各模块接口暴露文档补入。

\textbf{设计原则}：
\begin{itemize}[nosep]
  \item 每条参数有稳定 \texttt{id}、建议工程名、单位、归属节号、写者、读者、状态；
  \item L0 禁止业务模块私建同义 magic number；
  \item L1 禁止 PHY 子模块自行改写（如自主升 MCS、改 $P_o$）；
  \item 文档会迭代：\texttt{frozen} / \texttt{draft} / \texttt{tbd} 标明置信度，不要求等齐全员终稿。
\end{itemize}

\section{状态与读写约定}

\begin{longtable}{L{2.2cm}L{12.5cm}}
\caption{状态定义\label{tab:status}} \\
\toprule
\textbf{status} & \textbf{含义} \\
\midrule
\texttt{frozen} & 当前工程可直接引用；数值来自已有分册且与标准一致（仍以正式标准为准） \\
\texttt{draft} & 已从本工作区文档抽出；字段名/结构可能随他章对齐而调整 \\
\texttt{tbd} & 已知需要，但本工作区缺少完整定义，待与对应负责人边界对齐 \\
\bottomrule
\end{longtable}

\begin{longtable}{L{2.2cm}L{12.5cm}}
\caption{读写约定\label{tab:rw}} \\
\toprule
\textbf{层} & \textbf{约定} \\
\midrule
L0 & 只读；由 \texttt{slb\_phy\_const}（或等价）统一导出 \\
L1 & 上层写、PHY 读；建议落入单一 \texttt{SlbPhySessionConfig}（定位子集可映到 \texttt{SlbPosContext}） \\
\bottomrule
\end{longtable}

\section{抽取来源}

\begin{longtable}{L{4.5cm}L{10.2cm}}
\caption{本册参数主要来源\label{tab:src}} \\
\toprule
\textbf{来源} & \textbf{贡献} \\
\midrule
6.2.2 技术文档 & $f_s$、$T_s$、$\Delta f$、符号/CP、超帧、TTI、161/DC、$N_{\mathrm{CH}}$/$L_{\mathrm{CH}}$ \\
6.5 章级/分册 & $P_o$、$P_{\mathrm{MAX}}$、参考功率、MCS$\rightarrow$调制阶、$N_{\mathrm{ID}}$ 类消费 \\
6.6 及接口暴露 & PhyID、角色、能力位、坐标系、XRC 已解析字段子集、光速 \\
第 7 / 8 章（引用） & XRC/能力 IE 归属、射频 $P_{\mathrm{MAX}}$（本册标 draft/tbd） \\
\bottomrule
\end{longtable}

\newpage
%======================================================================
\part{L0 系统常量}
%======================================================================

\section{L0 说明}

L0 是全物理层共用的只读真源。任何 6.2 / 6.5 / 6.6 实现库\textbf{不得}在本地再写一份 ``120\,kHz'' ``161'' 等字面量副本；应引用本表 \texttt{id} 或生成头文件中的常量名。

\section{时间与频率基准}

\begin{longtable}{L{3.6cm}L{2.8cm}L{1.6cm}L{2.2cm}L{1.6cm}L{1.4cm}L{1.6cm}}
\caption{L0：时间与频率基准\label{tab:l0-tf}} \\
\toprule
\textbf{id} & \textbf{工程名} & \textbf{符号} & \textbf{取值} & \textbf{单位} & \textbf{归属} & \textbf{状态} \\
\midrule
\endfirsthead
\multicolumn{7}{c}{\small（续表）时间与频率基准} \\
\toprule
\textbf{id} & \textbf{工程名} & \textbf{符号} & \textbf{取值} & \textbf{单位} & \textbf{归属} & \textbf{状态} \\
\midrule
\endhead
\bottomrule
\endfoot
\texttt{slb.phy.fsHz} &
\texttt{fsHz} &
$f_s$ &
$30.72\times10^{6}$ &
Hz &
6.1 &
\texttt{frozen} \\
\texttt{slb.phy.tsSec} &
\texttt{tsSec} &
$T_s$ &
$1/f_s\approx 32.55\times10^{-9}$ &
s &
6.1 &
\texttt{frozen} \\
\texttt{slb.phy.deltaFHz} &
\texttt{deltaFHz} &
$\Delta f$ &
$120\times10^{3}$（$=f_s/256$） &
Hz &
6.2.2.1 &
\texttt{frozen} \\
\texttt{slb.phy.ofdmUsefulLenTs} &
\texttt{ofdmUsefulLenTs} &
--- &
$256$ &
$T_s$ &
6.2.2.1 &
\texttt{frozen} \\
\texttt{slb.phy.ofdmUsefulDurationUs} &
\texttt{ofdmUsefulDurationUs} &
--- &
$\approx 8.333$ &
$\mu$s &
6.2.2.1 &
\texttt{frozen} \\
\texttt{slb.phy.superframeDurationUs} &
\texttt{superframeDurationUs} &
$T_f$ &
$1000$（$=30720\,T_s$） &
$\mu$s &
6.2.2.2 &
\texttt{frozen} \\
\texttt{slb.phy.radioFramesPerSuperframe} &
\texttt{radioFramesPerSuperframe} &
--- &
$8$ &
--- &
6.2.2.2 &
\texttt{frozen} \\
\end{longtable}

\begin{longtable}{L{3.8cm}L{3.2cm}L{8.2cm}}
\caption{L0 时频基准：读写\label{tab:l0-tf-rw}} \\
\toprule
\textbf{id} & \textbf{写} & \textbf{读（消费者）} \\
\midrule
上表全部 & 无（常量） & 6.2.2--6.2.10、6.5.*、6.6.*、基带定时 \\
\bottomrule
\end{longtable}

\section{资源网格与载波常量}

\begin{longtable}{L{3.8cm}L{3.0cm}L{2.4cm}L{1.6cm}L{1.8cm}L{1.6cm}}
\caption{L0：网格与载波常量\label{tab:l0-grid}} \\
\toprule
\textbf{id} & \textbf{工程名} & \textbf{取值} & \textbf{单位} & \textbf{归属} & \textbf{状态} \\
\midrule
\endfirsthead
\multicolumn{6}{c}{\small（续表）网格与载波} \\
\toprule
\textbf{id} & \textbf{工程名} & \textbf{取值} & \textbf{单位} & \textbf{归属} & \textbf{状态} \\
\midrule
\endhead
\bottomrule
\endfoot
\texttt{slb.phy.subcarrierCount} &
\texttt{subcarrierCount} &
$161$ &
--- &
6.2.2.4 &
\texttt{frozen} \\
\texttt{slb.phy.dcSubcarrierIndex} &
\texttt{dcSubcarrierIndex} &
$80$（不映射） &
--- &
6.2.2.4 / 6.5.3 &
\texttt{frozen} \\
\texttt{slb.phy.baseChannelBandwidthMHz} &
\texttt{baseChannelBandwidthMHz} &
$20$ &
MHz &
6.2.2.4 &
\texttt{frozen} \\
\texttt{slb.phy.baseSubcarrierGroupCount} &
\texttt{baseSubcarrierGroupCount} &
$16$（每 10 有效子载波一组） &
--- &
6.2.2.4 &
\texttt{frozen} \\
\texttt{slb.phy.rsQpskLength} &
\texttt{rsQpskLength} &
$161$ &
--- &
6.5.3 &
\texttt{frozen} \\
\texttt{slb.phy.goldPnForRs} &
\texttt{goldPnLengthForRs} &
$M_{\mathrm{PN}}=322$ &
bit &
6.5.3 &
\texttt{frozen} \\
\end{longtable}

\noindent\textbf{读写}：写=无；读=资源映射、RS 生成、SC-FDMA/OFDM、定位 PMS 映射等。

\section{符号类型相关常量（查表用）}

以下为 L0\textbf{查表常量}（随符号类型取一行），表体来自 6.2.2.1 / 6.2.2.2；运行时由 L1 的符号类型配置选择行。

\begin{longtable}{L{2.0cm}L{2.4cm}L{2.4cm}L{2.0cm}L{2.4cm}C{1.6cm}}
\caption{L0：五种符号类型 CP/时长摘要\label{tab:l0-sym}} \\
\toprule
\textbf{类型} & \textbf{一般 CP} & \textbf{一般符号} & \textbf{GAP1} & \textbf{边界符号} & \textbf{$N_{\mathrm{sym}}^{\mathrm{frame}}$} \\
\midrule
\endfirsthead
\multicolumn{6}{c}{\small（续表）符号类型} \\
\toprule
\textbf{类型} & \textbf{一般 CP} & \textbf{一般符号} & \textbf{GAP1} & \textbf{边界符号} & \textbf{$N_{\mathrm{sym}}^{\mathrm{frame}}$} \\
\midrule
\endhead
\bottomrule
\endfoot
一 A & $18T_s$ & $274T_s$ & $20T_s$ & $276T_s$ & 14 \\
一 B & $18T_s$ & $274T_s$ & $22T_s$ & $278T_s$ & 14 \\
二 & $39T_s$ & $295T_s$ & $44T_s$ & $300T_s$ & 13 \\
三 & $64T_s$ & $320T_s$ & --- & --- & 12 \\
四 & $128T_s$ & $384T_s$ & --- & --- & 10 \\
\end{longtable}

\begin{longtable}{L{4.5cm}L{10.2cm}}
\caption{L0：符号类型与 $N_{\mathrm{CP}}$ 映射（供 6.5.3）\label{tab:l0-ncp}} \\
\toprule
\textbf{id / 规则} & \textbf{内容} \\
\midrule
\texttt{slb.phy.symbolTypeToNcp} &
类型一 A/B $\rightarrow N_{\mathrm{CP}}=3$；类型二 $\rightarrow 2$；类型三 $\rightarrow 1$；类型四 $\rightarrow 0$ \\
状态 & \texttt{frozen}；归属 6.2.2 $\leftrightarrow$ 6.5.3 \\
写 / 读 & 无 / RS QPSK、加扰相关初始化 \\
\bottomrule
\end{longtable}

\section{物理常数}

\begin{longtable}{L{3.8cm}L{3.2cm}L{3.5cm}L{1.6cm}L{1.6cm}L{1.6cm}}
\caption{L0：物理常数\label{tab:l0-phys}} \\
\toprule
\textbf{id} & \textbf{工程名} & \textbf{取值} & \textbf{单位} & \textbf{归属} & \textbf{状态} \\
\midrule
\texttt{slb.phy.speedOfLightMps} &
\texttt{speedOfLightMps} &
$299792458.0$（默认） &
m/s &
6.6.2.1 &
\texttt{frozen} \\
\end{longtable}

\noindent\textbf{读写}：写=无（测试可注入覆盖，但不得改默认真源）；读=RTT/复合 CSI 解距等。

\newpage
%======================================================================
\part{L1 会话 / 上层下发配置}
%======================================================================

\section{L1 说明}

L1 是\textbf{会话级}配置：在接入、XRC 重配、调度更新时由上层写入，PHY 各库通过只读上下文消费。推荐工程入口：

\begin{itemize}[nosep]
  \item 全 PHY：\texttt{SlbPhySessionConfig}（待下一步固化字段表）；
  \item 定位子集：已有 \texttt{SlbPosContext} / \texttt{slbPosApplyXrcConfig}（6.6 接口暴露），应映射到同一 L1 真源，避免双写。
\end{itemize}

\section{身份与节点}

\begin{longtable}{L{3.4cm}L{2.6cm}L{1.4cm}L{2.0cm}L{2.4cm}L{2.4cm}L{1.4cm}}
\caption{L1：身份与节点\label{tab:l1-id}} \\
\toprule
\textbf{id} & \textbf{工程名} & \textbf{单位} & \textbf{归属} & \textbf{写} & \textbf{读} & \textbf{状态} \\
\midrule
\endfirsthead
\multicolumn{7}{c}{\small（续表）身份与节点} \\
\toprule
\textbf{id} & \textbf{工程名} & \textbf{单位} & \textbf{归属} & \textbf{写} & \textbf{读} & \textbf{状态} \\
\midrule
\endhead
\bottomrule
\endfoot
\texttt{slb.session.localPhyId} &
\texttt{localPhyId} &
--- &
接入/XRC &
第 7 章/MAC &
6.6 FLAG/报告 &
\texttt{draft} \\
\texttt{slb.session.cellId} &
\texttt{nId} / $N_{\mathrm{ID}}$ &
--- &
6.5.3/小区 &
系统消息/接入 &
RS、部分 $c_{\mathrm{init}}$ &
\texttt{draft} \\
\texttt{slb.session.gPid} &
\texttt{nGPid} / $N_{\mathrm{G-PID}}$ &
--- &
6.5.7 &
调度/系统 &
数据加扰默认 $c_{\mathrm{init}}$ &
\texttt{draft} \\
\texttt{slb.session.posRole} &
\texttt{role} &
枚举 &
6.6.1 &
XRC &
6.6 全包 &
\texttt{draft} \\
\texttt{slb.session.fixNodePhyId} &
\texttt{fixNodePhyId} &
--- &
6.6.1 &
XRC &
6.6.4 解算 &
\texttt{draft} \\
\end{longtable}

\section{帧、TTI 与载波配置}

\begin{longtable}{L{3.6cm}L{2.8cm}L{1.4cm}L{2.0cm}L{2.2cm}L{2.2cm}L{1.4cm}}
\caption{L1：帧 / TTI / 载波\label{tab:l1-frame}} \\
\toprule
\textbf{id} & \textbf{工程名} & \textbf{单位} & \textbf{归属} & \textbf{写} & \textbf{读} & \textbf{状态} \\
\midrule
\endfirsthead
\multicolumn{7}{c}{\small（续表）帧/TTI/载波} \\
\toprule
\textbf{id} & \textbf{工程名} & \textbf{单位} & \textbf{归属} & \textbf{写} & \textbf{读} & \textbf{状态} \\
\midrule
\endhead
\bottomrule
\endfoot
\texttt{slb.session.symbolType} &
\texttt{symbolType} &
枚举 &
6.2.2.1 &
高层/MAC &
帧结构、OFDM、$N_{\mathrm{CP}}$ &
\texttt{draft} \\
\texttt{slb.session.nTti} &
\texttt{nTti} &
0--6 &
6.2.2.3 &
高层/MAC &
TTI 长度=$2^{N_{\mathrm{TTI}}}$ 无线帧 &
\texttt{draft} \\
\texttt{slb.session.nCh} &
\texttt{nCh} &
档位 &
6.2.2.4 &
高层/MAC &
工作信道数 &
\texttt{draft} \\
\texttt{slb.session.lCh} &
\texttt{lCh} &
档位 &
6.2.2.4 &
由 $N_{\mathrm{CH}}$ 表定或配置 &
工作载波粒度 &
\texttt{draft} \\
\texttt{slb.session.centerFreqHz} &
\texttt{centerFreqHz} &
Hz &
8.3 &
射频/系统 &
多载波布局 &
\texttt{tbd} \\
\end{longtable}

\noindent\textbf{注}：$N_{\mathrm{CH}}$ 与 $L_{\mathrm{CH}}$ 18 档对应关系见 6.2.2 表（本册不重复抄写全部档位）；实现时应引用同一查表，禁止模块内硬编码残缺子集。

\section{功率相关（上层下发）}

\begin{longtable}{L{3.8cm}L{3.2cm}L{1.4cm}L{2.0cm}L{2.4cm}L{2.0cm}L{1.4cm}}
\caption{L1：功率配置\label{tab:l1-pwr}} \\
\toprule
\textbf{id} & \textbf{工程名} & \textbf{单位} & \textbf{归属} & \textbf{写} & \textbf{读} & \textbf{状态} \\
\midrule
\endfirsthead
\multicolumn{7}{c}{\small（续表）功率} \\
\toprule
\textbf{id} & \textbf{工程名} & \textbf{单位} & \textbf{归属} & \textbf{写} & \textbf{读} & \textbf{状态} \\
\midrule
\endhead
\bottomrule
\endfoot
\texttt{slb.session.p0SrsDbm} &
\texttt{srsP0NominalDbm} &
dBm &
6.5.4 &
XRC/系统消息 &
6.5.4 &
\texttt{draft} \\
\texttt{slb.session.p0RachDbm} &
\texttt{rachP0NominalDbm} &
dBm &
6.5.4 &
同上 &
6.5.4 &
\texttt{draft} \\
\texttt{slb.session.p0AckDbm} &
\texttt{ackP0NominalDbm} &
dBm &
6.5.4 &
同上 &
6.5.4 &
\texttt{draft} \\
\texttt{slb.session.p0DataClass1Dbm} &
\texttt{dataClass1P0NominalDbm} &
dBm &
6.5.4 &
同上 &
6.5.4 &
\texttt{draft} \\
\texttt{slb.session.p0DataClass2Dbm} &
\texttt{dataClass2P0NominalDbm} &
dBm &
6.5.4 &
同上 &
6.5.4 &
\texttt{draft} \\
\texttt{slb.session.gLinkRefSignalPowerDbm} &
\texttt{gLinkReferenceSignalPowerDbm} &
dBm &
6.5.4 &
XRC/系统消息 &
$P_L$、功控 &
\texttt{draft} \\
\texttt{slb.session.pMaxDbm} &
\texttt{pMaxDbm} &
dBm &
第 8 章/射频 &
射频能力/法规 &
6.5.4 限幅 &
\texttt{draft} \\
\end{longtable}

\noindent\textbf{强制}：PHY 只执行已配置的 $P_o$/$P_{\mathrm{MAX}}$，不得根据信道质量自行改写目标功率。

\section{MCS 与调制阶}

\begin{longtable}{L{3.6cm}L{2.8cm}L{1.4cm}L{2.0cm}L{2.4cm}L{2.2cm}L{1.4cm}}
\caption{L1：MCS / 调制\label{tab:l1-mcs}} \\
\toprule
\textbf{id} & \textbf{工程名} & \textbf{单位} & \textbf{归属} & \textbf{写} & \textbf{读} & \textbf{状态} \\
\midrule
\endfirsthead
\multicolumn{7}{c}{\small（续表）MCS} \\
\toprule
\textbf{id} & \textbf{工程名} & \textbf{单位} & \textbf{归属} & \textbf{写} & \textbf{读} & \textbf{状态} \\
\midrule
\endhead
\bottomrule
\endfoot
\texttt{slb.session.mcsIndex} &
\texttt{mcsIndex} &
--- &
6.2.5 &
调度/MAC &
查 MCS 表 $\rightarrow$ 调制阶 &
\texttt{draft} \\
\texttt{slb.session.modType} &
\texttt{modType} &
枚举 &
6.5.6 &
由 MCS/固定配置导出 &
6.5.6 调制 LUT &
\texttt{draft} \\
\texttt{slb.session.modulationOrderM} &
\texttt{modulationOrderM} &
$\{2,4,\ldots,12\}$ &
6.5.6 &
派生 &
比特打包/符号数 &
\texttt{draft} \\
\end{longtable}

\noindent\textbf{链路特例（配置面，非 PHY 自主）}：多数控制固定 QPSK；TCI 二阶段 ACK 可为 QPSK/16QAM；数据为 MCS 表决定的 QPSK--4096QAM。\\
\textbf{强制}：\texttt{【无自主升阶】}——PHY 不根据 SINR 自行改 MCS。

\section{定位能力与坐标系}

\begin{longtable}{L{3.8cm}L{3.0cm}L{1.4cm}L{2.0cm}L{2.2cm}L{2.2cm}L{1.4cm}}
\caption{L1：定位能力 / 坐标系\label{tab:l1-poscap}} \\
\toprule
\textbf{id} & \textbf{工程名} & \textbf{单位} & \textbf{归属} & \textbf{写} & \textbf{读} & \textbf{状态} \\
\midrule
\endfirsthead
\multicolumn{7}{c}{\small（续表）定位能力} \\
\toprule
\textbf{id} & \textbf{工程名} & \textbf{单位} & \textbf{归属} & \textbf{写} & \textbf{读} & \textbf{状态} \\
\midrule
\endhead
\bottomrule
\endfoot
\texttt{slb.session.ofdmPositioningCapable} &
\texttt{ofdmPositioningCapable} &
bool &
6.6 / 能力 IE &
第 7 章能力 &
6.6 RTT 准入 &
\texttt{draft} \\
\texttt{slb.session.posFrameType} &
\texttt{frameType} &
ABS/REL &
6.6.1/6.6.4 &
XRC &
坐标解算输出 &
\texttt{draft} \\
\texttt{slb.session.anchorXyzM} &
\texttt{anchorXyzM[3]} &
m &
6.6.1 &
XRC &
6.6.4 &
\texttt{draft} \\
\texttt{slb.session.syncSequenceU} &
\texttt{syncSequenceU} &
--- &
同步/FLAG &
G 配置 &
与 PMS $u$ 隔离校验 &
\texttt{draft} \\
\end{longtable}

\section{XRC 已解析字段（定位配置子集）}

本表只登记\textbf{PHY/6.6 允许消费的已解析字段}；禁止传入 ASN.1 原文。完整 IE 以第 7.5 章为准（多数字段 \texttt{tbd}/\texttt{draft}）。

\begin{longtable}{L{3.6cm}L{3.2cm}L{5.0cm}L{2.2cm}L{1.4cm}}
\caption{L1：XRC 已解析字段（子集）\label{tab:l1-xrc}} \\
\toprule
\textbf{id} & \textbf{工程名/字段} & \textbf{用途} & \textbf{读} & \textbf{状态} \\
\midrule
\endfirsthead
\multicolumn{5}{c}{\small（续表）XRC 字段} \\
\toprule
\textbf{id} & \textbf{工程名/字段} & \textbf{用途} & \textbf{读} & \textbf{状态} \\
\midrule
\endhead
\bottomrule
\endfoot
\texttt{slb.session.xrc.roles} &
角色指派 / \texttt{posCalculator} 等 &
锚点/标签/解算节点 &
6.6.1 &
\texttt{draft} \\
\texttt{slb.session.xrc.staticRes} &
静态时频/端口/波束/载波/带宽/序列 &
纯 XRC 模式全量；或与 GCI 组合 &
6.6.* &
\texttt{draft} \\
\texttt{slb.session.xrc.measItems} &
可选测量项开关 &
RTT/角度/跳频等任务使能 &
6.6 &
\texttt{draft} \\
\texttt{slb.session.xrc.groupcast} &
\texttt{groupcastConfig} &
FLAG 组播 &
6.6.2.2 &
\texttt{draft} \\
\texttt{slb.session.xrc.nonConnectMeas} &
\texttt{nonConnectMeasConfig} &
无连接测量 &
6.6 &
\texttt{draft} \\
\texttt{slb.session.xrc.nodeCoords} &
\texttt{nodePositioningCoordinates} &
锚点已知坐标 &
6.6.4 &
\texttt{draft} \\
\texttt{slb.session.xrc.reportTarget} &
上报目标 &
报告 IE 汇聚方向 &
6.6 P7 &
\texttt{draft} \\
\texttt{slb.session.xrc.hopConfig} &
跳频图案 / $T_{\mathrm{RF}}$ 等 &
6.6.2.3 复合测距 &
6.6 P4 &
\texttt{draft} \\
\texttt{slb.session.xrc.securePms} &
安全 PMS 开关 &
是否用安全序列 &
6.2.3.9 / 6.6 &
\texttt{draft} \\
\end{longtable}

\noindent\textbf{写}：一律为第 7 章解析器 $\rightarrow$ \texttt{slbPosApplyXrcConfig}（或将来的 SessionConfig 应用接口）。\\
\textbf{禁止}：6.6/PHY 内二次 ASN.1 解码；禁止暴露射频句柄与内部缓冲布局。

\section{与动态指示的边界（非 L1 常驻，但相关）}

下列量常由 GCI/TCI \textbf{当次}更新，更接近 L3 事件入参；此处仅作边界说明，避免误放入 Session 常驻字段：

\begin{longtable}{L{4.5cm}L{10.2cm}}
\caption{L1 与当次动态指示的边界\label{tab:l1-vs-dyn}} \\
\toprule
\textbf{对象} & \textbf{处理建议} \\
\midrule
GCI 格式 0--2 资源 & 当次 RTT/角度时频；经 \texttt{slbPosApplyGciResource} 写入会话的``当前资源''槽，勿与 XRC 静态混为同一不可变字段 \\
GCI 第四格式激活位 & 事件：开测/关测；不作为长期配置常量 \\
调度给出的当次 $M$（子载波数） & 功控入参；可视为 L3/调用参数，不必全部进入 Session 静态区 \\
\bottomrule
\end{longtable}

\newpage
%======================================================================
\part{汇总、缺口与下一步}
%======================================================================

\section{L0 / L1 条目计数（初稿）}

\begin{longtable}{L{3cm}C{3cm}L{8.5cm}}
\caption{初稿覆盖统计\label{tab:count}} \\
\toprule
\textbf{层} & \textbf{约条目} & \textbf{说明} \\
\midrule
L0 & $\sim$20+（含符号类型表） & 时频、网格、符号查表、$N_{\mathrm{CP}}$、光速 \\
L1 身份/帧/载波 & $\sim$10 & PhyID、小区/G-PID、角色、$N_{\mathrm{TTI}}$、$N_{\mathrm{CH}}$ 等 \\
L1 功率/MCS & $\sim$10 & 五类 $P_o$、参考功率、$P_{\mathrm{MAX}}$、MCS/调制 \\
L1 定位/XRC & $\sim$12 & 能力、坐标、XRC 已解析子集 \\
\bottomrule
\end{longtable}

\section{已知缺口（需向他人/他章对齐）}

\begin{longtable}{L{4.5cm}L{10.2cm}}
\caption{tbd / 高优先级 draft\label{tab:gap}} \\
\toprule
\textbf{缺口} & \textbf{建议动作} \\
\midrule
完整 XRC IE 字段表 & 向第 7.5 章负责人只要「已解析结构体字段清单」，不要等全文 \\
$N_{\mathrm{ID}}$ / $N_{\mathrm{G-PID}}$ 正式工程名与位宽 & 与系统消息/接入文档对齐后改 \texttt{frozen} \\
中心频率与信道号表 & 对齐第 8.3 章 \\
MCS 表索引 $\leftrightarrow$ 调制阶完整映射 & 对齐 6.2.5 分册后写入查表（仍属 L1 派生） \\
SessionConfig 与 PosContext 字段合一 & 下一步专项：避免双会话 \\
\bottomrule
\end{longtable}

\section{建议的下一步工作}

\begin{enumerate}[nosep]
  \item 组内评审本 PDF：确认 L0 可全部 \texttt{frozen}；
  \item 起草 \texttt{SlbPhySessionConfig} 字段表（直接映射本册 L1）；
  \item 编写 \texttt{slb-param-contract} Skill（强制引用本登记册）；
  \item 按模块补 L2，并推广「接口暴露」文档模式；
  \item （可选）导出 \texttt{registry.yaml} 供代码生成与 magic-number 校验。
\end{enumerate}

\section{约束摘要（实现取向）}

\begin{enumerate}[nosep]
  \item \textbf{【无状态机】}：配置应用与流程推进用事件/条件判断，不维护章级 FSM；
  \item \textbf{【无自主升阶】}：MCS/调制阶只执行配置；
  \item \textbf{【XRC 已解析】}：禁止 raw PDU 进入 PHY；
  \item \textbf{【单位后缀】}：工程名遵循 nearlink-naming（\texttt{Hz}/\texttt{Us}/\texttt{Dbm}/\texttt{M} 等）；
  \item \textbf{【分册为准】}：与正式标准或更新分册冲突时，先改登记册再改代码。
\end{enumerate}

\vspace{1.5em}
{\small\textcolor{gray}{本登记册为工程/AI 实现用初稿，仅供理解与对齐参考；精确参数与字段定义以 T/XS 10001-2025 正式标准及各有效技术文档分册为准。}}

\end{document}
